\begin{statement}{C.2.UN.58}{Proved}{The complete canonical and period reduction for all-Hopf trees}
Let $S$ be a full reduced SNC fibre of a smooth-total degeneration
of connected complex two-tori, with every irreducible component a
primary or secondary Hopf surface. If one component is a diagonal
Hopf surface of algebraic dimension zero, the whole support is
one of the diagonal cycles classified by C.2.UN.46 and C.2.UN.48.

In every remaining case the incidence graph is the unicyclic graph
of C.2.UN.54. Every cycle vertex and every nonterminal tree vertex
has algebraic dimension one. Its actual elliptic reduction is
\[
                    q_v:H_v\longrightarrow\PP^1,        \tag{C.HT.1}
\]
with fixed Jacobian, identity integral monodromy and smooth multiple
supports. The other possible vertices are nonlinear
algebraic-dimension-zero Hopf leaves. Every multiple fibre of
(C.HT.1) is a double curve of $S$.

For each edge $e$ the normal line has a finite exact order $m_e$.
It is the full fibre multiplicity on both elliptic Hopf endpoints,
and is the resonance at a nonlinear endpoint. The connected
normal-root covers of the two endpoints are identified by the
actual dual-normal map. Their Jacobians therefore identify along
every edge, and the product of their linear identifications around
the cycle is the identity. The actual elliptic cycle product can
still be a translation, whose value is retained by these maps.

Let $k_v$ be the canonical evaluation orders, and let $k_0$ be their
common minimum on the cycle. Root each attached tree at that cycle.
The heights are determined by the weighted graph:
\[
 h_v:=k_v-k_0=
 \begin{cases}
 0,&v\text{ is on the cycle},\\
 \displaystyle\sum_{e\text{ on the path from the cycle to }v}m_e,
       &v\text{ is in an attached tree}.
 \end{cases}                                           \tag{C.HT.2}
\]
More precisely, after its invariant relative elliptic one-form is
factored out, the canonical form on a cycle vertex is a nonzero
constant multiple of $dz/z$, where the two cycle points are
$0,\infty$. On a tree vertex it is a nonzero constant multiple of
$dz$, where its parent point is $\infty$.

All regular elliptic-fibre two-tori on the algebraic-dimension-one
vertices, together with the real core two-tori on the nonlinear
leaves, represent one transported integral class $\gamma$ in the
nearby fibre. If a nonlinear leaf occurs, the single period
\[
                     P(t)=\int_{\gamma_t}\omega_t       \tag{C.HT.3}
\]
of the relative canonical evaluation divided by $t^{k_0}$ has a
finite positive vanishing order $R$. Every nonlinear leaf has
height $R$, and every algebraic-dimension-one vertex has height
strictly smaller than $R$. These are assertions about the actual
common period. When there is no nonlinear leaf its period is
allowed to vanish identically.
\end{statement}

\begin{proof}
\par\medskip\noindent\textbf{Separation by the actual normal lines.}
The axes of an algebraic-dimension-zero diagonal primary Hopf
surface have nontorsion normal lines. For example, torsion of the
normal of $w=0$ for $(z,w)\mapsto(\alpha z,\beta w)$ would give
$\beta^n=\alpha^a$ for some integers $n>0,a$. The corresponding
monomial quotient is a nonconstant meromorphic invariant, contrary
to algebraic dimension zero. This is the exact automorphy
calculation of C.2.UN.38. Pullback to the finite primary cover
proves nontorsion on a secondary quotient as well.

Every curve on an algebraic-dimension-one Hopf surface is a fibre
support of (C.HT.1): a horizontal curve would meet a regular elliptic
fibre, contradicting $H^2(H,\Q)=0$ and positivity of intersection.
Its normal has finite order, as checked explicitly below. The
nonlinear Hopf axis also has finite order equal to its resonance,
by (C.SN.3)--(C.SN.4). Normal duality at an actual reduced SNC edge
preserves torsion and nontorsion. Thus a diagonal
algebraic-dimension-zero vertex can meet only another such vertex.
Connectedness then puts every component in that class, and
C.2.UN.54 gives precisely the already classified diagonal cycle.

In the remaining case every algebraic-dimension-zero vertex is
nonlinear and has just its one elliptic curve. It is a leaf.
All other vertices have (C.HT.1) by C.2.UN.37. In particular every
cycle or nonterminal tree vertex has that fibration.

\par\medskip\noindent\textbf{The root charts and the invariant relative form.}
At a full fibre $mJ$ of (C.HT.1), its actual normalized $m$th root
cover is a connected smooth elliptic family, finite unramified over
the original surface tube. The fixed modulus and identity linear
monodromy give a local principal bundle with fixed elliptic fibre
$F$. It is holomorphically trivial on a disc: the exponential
sequence $0\to\Lambda_F\to\OO\to F\to0$, together with
$H^1(\Delta,\OO)=H^2(\Delta,\Lambda_F)=0$, kills its translation
torsor. Thus the original tube is
\[
       (F\times\Delta_s)/\mu_m,\qquad
       (z,s)\longmapsto(z+\tau,\zeta_m s),\qquad
       q=s^m,\quad \operatorname{ord}(\tau)=m.          \tag{C.HT.4}
\]
Here the translation can be made constant by an actual gauge.
If its initially holomorphic expression is $\tau+A(s)$, choose a
small complex lift with $\sum_{j=0}^{m-1}A(\zeta_m^j s)=0$ and put
$b(s)=-m^{-1}\sum_{j=0}^{m-1}jA(\zeta_m^j s)$.
Then $b(\zeta_ms)-b(s)=-A(s)$, so $z\mapsto z+b(s)$ gives (C.HT.4).
The exact order follows from the free deck action on the connected
central elliptic curve.

The normal of $J=F/\langle\tau\rangle$ is the associated line for
the primitive character $\zeta_m$. It has exact order $m$:
if a smaller tensor power were trivial, its pullback trivialization
would be a nowhere-zero holomorphic function on the compact $F$
with a nontrivial character, which is impossible. The connected
normal-root cover is exactly $F\to J$. Duality of the two normal
lines at an edge identifies their exact orders and their connected
root covers. Hence the number $m_e$ in the statement is well defined,
including at a nonlinear leaf via the exact secondary charts of
C.2.UN.50.

In every root chart write $dz$ for a fixed nonzero invariant
elliptic one-form. Transition maps between principal elliptic
charts are translations; their differentials differ by a multiple
of the base differential. Consequently
\[
                         q^*\theta\wedge dz             \tag{C.HT.5}
\]
is a globally defined meromorphic two-form whenever $\theta$ is a
meromorphic differential on $\PP^1$. The finite translations in
(C.HT.4) preserve it. This constructs the needed form rather than
postulating a global absolute elliptic one-form.

Every meromorphic two-form whose divisor is vertical has precisely
the form (C.HT.5). Its ratio to one such reference form is a
meromorphic function with vertical divisor. On a regular compact
elliptic fibre away from that divisor the ratio is a nonzero
holomorphic function and is constant. It therefore descends on
the regular base. In (C.HT.4) its descent is a meromorphic function
of $s$ invariant under $s\mapsto\zeta_ms$, and hence is a
meromorphic function of $s^m$. This proves descent across every
multiple fibre as well. In this identification the exact order is
\[
 \operatorname{ord}_J(q^*\theta\wedge dz)
               =m\,\operatorname{ord}_{q(J)}\theta+m-1 .
                                                        \tag{C.HT.6}
\]

\par\medskip\noindent\textbf{The canonical equation and the leaf induction.}
The actual canonical evaluation gives on every component a
nonzero meromorphic two-form $\eta_v$ with divisor
\[
 \operatorname{div}(\eta_v)
       =\sum_{e=(v,w)}(k_w-k_v-1)J_e .                 \tag{C.HT.7}
\]
This is adjunction with the specified vertical orders. It can also
be read in the double chart $t=xy$ by dividing the evaluation by
$t^{k_v}$ and taking the residue after division by $t$. Away from
the boundary there is no zero or pole, since the evaluation has
only its central component divisors. Equation (C.HT.7) is a divisor
of the actual section, not just numerical equivalence.

Apply (C.HT.6) at an algebraic-dimension-one vertex. An unmarked
multiple fibre would have order zero in (C.HT.7), requiring
$m-1\equiv0\pmod m$ and hence $m=1$. There are therefore no
unmarked multiple fibres. At a marked fibre (C.HT.6) says
\[
 d_{v,e}:=\frac{k_w-k_v}{m_e}
       =\operatorname{ord}_{q_v(J_e)}\theta_v+1\in\Z.
                                                        \tag{C.HT.8}
\]
Every other point of $\PP^1$ has order zero. Taking the degree of
the divisor of this rational differential proves
\[
                    \sum_{e\ni v}d_{v,e}
                           =\operatorname{val}(v)-2.   \tag{C.HT.9}
\]
Opposite orientations give opposite $d$ because the same exact
normal order $m_e$ occurs at the two endpoints.

C.2.UN.54 identifies the cycle vertices exactly as the minimum
canonical-order vertices. Their $k_v$ are one common $k_0$,
so both cycle-edge differences are zero. A terminal
algebraic-dimension-one vertex has just its parent edge, and
(C.HT.9) gives $k_{\rm leaf}-k_{\rm parent}=m_e$.
For a nonlinear leaf the invariant form
$dz\wedge dw/w^{r+1}$ and the infinite order of its effective
axis line give exactly the same equality, with $m_e=r$,
as in (C.AS.12).

Now work inward along any finite attached tree. Suppose all the
$c$ child edges of a vertex have outgoing $d=1$. At that vertex
the one parent term is the negative of the parent's outgoing
$d_{\rm parent}$, and (C.HT.9) gives
$-d_{\rm parent}+c=c-1$. Thus $d_{\rm parent}=1$.
Every nonterminal tree vertex has algebraic dimension one, so this
induction reaches the cycle. It proves (C.HT.2) on every edge.

At a cycle vertex the two cycle terms have $d=0$, giving simple
poles of $\theta_v$ there; every child term has $d=1$ and order
zero. There are no other zeros or poles. A coordinate taking the
cycle points to $0,\infty$ therefore gives $\theta_v=c\,dz/z$.
At a tree vertex the parent term has $d=-1$ and order $-2$;
all child terms have order zero. A coordinate with its parent
point at infinity gives $\theta_v=c\,dz$. In each case $c\ne0$.
This proves the more precise differential assertion.

\par\medskip\noindent\textbf{The actual Jacobian comparison around the cycle.}
The same root charts show $R^1q_{v*}\OO_{H_v}=\OO_{\PP^1}$.
On $F\times\Delta_s$ the cohomology frame is the fixed
$H^1(F,\OO_F)$; finite deck translations act trivially, so its
invariant functions are exactly functions of $s^m$. On overlaps
the identity linear monodromy and translations give the same
frame. Leray on $\PP^1$ therefore identifies
$H^1(H_v,\OO_{H_v})$ with this one-dimensional constant fibre
cohomology.

Restriction to $J_e$ followed by pullback along the actual finite
cover $F\to J_e$ agrees with evaluation on the regular elliptic
fibre. This follows from the commuting pullback diagram of the
surface tube and (C.HT.4), where both maps are restriction of the
same constant cohomology frame. Thus the scalar cycle product
in (C.HR.2) is exactly the action on $H^1(F,\OO_F)$ of the product
of the lifted fibre isomorphisms. It is one by C.2.UN.54.
An elliptic automorphism whose action on $H^1(\OO)$ is the identity
has linear part one: after choosing an origin, its lift is
$z\mapsto az+b$, and the nonzero invariant one-form detects $a$
(or its conjugate on $H^1(\OO)$). Hence that product is a
translation. Different choices of a lift through a finite
normal-root cover alter it only by deck translations. The
Jacobian product is therefore the identity with all maps retained.

\par\medskip\noindent\textbf{One integral core class and its nonzero nonlinear periods.}
We describe the cycle comparison before integrating. Near an
edge the two surface tubes are the disk bundles of the actual
torsion normal lines in (C.HT.4). Their dual-normal isomorphism
identifies their connected root covers. On that cover the
normal-product collar is
\[
                  xy=t,\qquad q_v=x^{m_e},\quad
                  q_w=y^{m_e}.                         \tag{C.HT.10}
\]
Multiplying one base coordinate by a nonzero constant makes these
equalities exact: the two trivializations of the $m_e$th normal
power differ by a holomorphic unit on a compact elliptic curve,
and hence by a constant. The regular fibre $x^{m_e}=c$ on the
first side is sent to $y^{m_e}=t^{m_e}/c$ on the other. Both are
the same connected elliptic cover of $J_e$, with its complex
orientation. Varying the nonzero normal level supplies the
oriented isotopy in the neck.

For an arbitrary given smoothing, the product-preserving smooth
tubular comparison used in C.2.UN.49 carries these normal collars
to the actual $xy=t$ collars. Compact cores are compared by
Ehresmann transport. Consequently (C.HT.10) is an equality of
the transported integral two-cycle classes in the actual nearby
fibre, without a holomorphic linearization of its neighbourhood.
Regular fibres on one component are isotopic through its regular
base, so every algebraic-dimension-one vertex gives the same
class $\gamma$. At a nonlinear leaf the exact unimodular
boundary maps (C.SN.6)--(C.SN.8), together with the displayed
smooth-pair deformation, identify its real core with this same
oriented class. They retain the full secondary order.

Write $\omega_t$ for the evaluation divided by $t^{k_0}$.
On a compact regular fibre of $H_v$ its leading coefficient is
$t^{h_v}\eta_v$. The restriction of the meromorphic two-form
$\eta_v$ to this holomorphic curve is zero. Therefore
\[
                     P(t)=O(t^{h_v+1})
       \quad\text{at every algebraic-dimension-one vertex}.
                                                        \tag{C.HT.11}
\]
These are actual periods of a closed relative holomorphic form
over the same compact class. They are holomorphic by
differentiation and Stokes' theorem; the local compact-core
transport gives the displayed extension and leading coefficient.

At a nonlinear leaf
$H_r(\beta)/\mu_n$ its unique-axis form is
$\eta=c\,dz\wedge dw/w^{r+1}$, $c\ne0$. In
$U=z/w^r$, $V=\log w$, the two primitive real core periods are
$(0,2\pi i/n)$ and $(\beta^{-r},\log\beta)$. Hence
\[
 P(t)=t^{h_{\rm leaf}}
       \left(\frac{\pm2\pi i c\beta^{-r}}{n}+O(t)\right),
                    \qquad c\ne0.                     \tag{C.HT.12}
\]
The coefficient is nonzero. One such leaf fixes the finite
vanishing order $R=h_{\rm leaf}>0$ of the single period.
Every other nonlinear leaf has that same height, while
(C.HT.11) gives $h_v<R$ at every elliptic Hopf vertex.
No nonvanishing conclusion is imposed when there is no nonlinear
leaf.

The exact root congruence, tree equations, height sums, collar
monomials and common-period inequalities are checked in
\path{code/E_hopf_tree_heights.py} and its saved output. The
argument proves these comparisons for every actual support in
the stated branch. It does not assert that these necessary
equations alone smooth arbitrary abstract attached-tree data.
\end{proof}


% C.2.INT.1 from part_C2_interfaces.tex

